% -----------------------------------------------------------------------------
% Chapter 22: Performance & Scalability
% -----------------------------------------------------------------------------
\chapter{Performance \& Scalability}
\label{chap:performance_scalability}

\section{Current Empirical Benchmark Profile}
Under the MVP release, CreditTech demonstrates exceptional execution performance on standard commodity hardware (single-node CPU execution on an 8-core workstation):
\begin{itemize}
  \item \textbf{Model Inference Latency:} The tuned Weight-of-Evidence (WoE) scorecard executes a complete scoring and SHAP attribution pass in \textbf{0.038\,ms} (over $26,000$ scoring evaluations per second per CPU core).
  \item \textbf{End-to-End API Response Time:} $P_{50}$ API latency across core ingestion and decision endpoints is \textbf{14.2\,ms}; $P_{95}$ latency remains below \textbf{48.0\,ms}.
  \item \textbf{Dataset Processing Throughput:} Ingested and processed all \textbf{22,000 benchmark records} (Home Credit + GMSC) through feature extraction and 5-Fold Cross-Validation in under \textbf{4.2 seconds}.
  \item \textbf{Concurrent Burst Resilience:} Successfully verified against a 60-request concurrent burst test (\texttt{tests/test\_p6\_load.py}) simulating 2$\times$ peak branch operational volume with 0 dropped requests.
\end{itemize}

\section{Scaling Roadmap: 100 to 100,000 Users}
Figure~\ref{fig:scaling_roadmap} illustrates the architectural evolution as the platform scales across four distinct user volume milestones.

\begin{figure}[H]
\centering
\begin{tikzpicture}[
  scale=0.92, every node/.style={transform shape},
  node distance=0.45cm, auto
]

  % Stage 1: MVP
  \node[clientblock, fill=primary!10, draw=primary, text width=3.35cm, minimum height=1.75cm] (s1) {
    \textbf{1. MVP Pilot}\\[1pt]
    \scriptsize 100 Active Borrowers\\
    \scriptsize Single-node SQLite\\
    \scriptsize In-memory caching
  };

  % Stage 2: Cluster Scale
  \node[serviceblock, right=0.35cm of s1, fill=accent!10, draw=accent, text width=3.35cm, minimum height=1.75cm] (s2) {
    \textbf{2. Regional Cluster}\\[1pt]
    \scriptsize 1,000 Active Users\\
    \scriptsize PostgreSQL DB Engine\\
    \scriptsize Redis Session Cache
  };

  % Stage 3: State-Wide
  \node[serviceblock, right=0.35cm of s2, fill=warn!10, draw=warn, text width=3.35cm, minimum height=1.75cm] (s3) {
    \textbf{3. State-Wide RRB}\\[1pt]
    \scriptsize 10,000 Active Users\\
    \scriptsize Celery / Redis Workers\\
    \scriptsize Read/Write DB Replicas
  };

  % Stage 4: National Scale
  \node[baseblock, right=0.35cm of s3, fill=success!15, draw=success, text width=3.35cm, minimum height=1.75cm] (s4) {
    \textbf{4. National Scale}\\[1pt]
    \scriptsize 100,000+ Borrowers\\
    \scriptsize Distributed Kubernetes\\
    \scriptsize Pre-computed EO tiles
  };

  \draw[flowline] (s1) -- (s2);
  \draw[flowline] (s2) -- (s3);
  \draw[flowline] (s3) -- (s4);
\end{tikzpicture}
\caption{Research Block Diagram: 4-Stage Architectural Scaling Roadmap from MVP to National Rural Deployment.}
\label{fig:scaling_roadmap}
\end{figure}

\section{Bottleneck Identification \& Scaling Mitigations}
Table~\ref{tab:bottleneck_analysis} details the primary technical bottlenecks identified during system stress profiling and the corresponding engineering mitigations.

\begin{table}[H]
\centering
\small
\renewcommand{\arraystretch}{1.2}
\begin{tabularx}{\textwidth}{p{3.2cm} p{3.4cm} X}
  \toprule
  \textbf{Subsystem} & \textbf{Potential Bottleneck} & \textbf{Architectural Mitigation} \\
  \midrule
  \textbf{Relational Database} & SQLite file-lock contention under concurrent writes. & Migrate to PostgreSQL with PgBouncer connection pooling and read-only replica routing for dashboard analytics. \\
  \addlinespace
  \textbf{Satellite Ingestion} & Third-party Copernicus API rate limits and GeoTIFF decoding latency. & Pre-compute 10m Sentinel-2 NDVI vegetative indices on a bi-weekly cron schedule at the \textit{Gram Panchayat} centroid level; store in Redis. \\
  \addlinespace
  \textbf{Account Aggregator} & Synchronous HTTP polling waiting on borrower mobile OTP approval. & Decouple via asynchronous webhook notifications: ingest service registers an event listener and pauses pipeline execution cleanly. \\
  \addlinespace
  \textbf{Background Jobs} & Python ASGI process thread contention during bulk batch scoring. & Offload heavy feature extraction and batch model scoring to Celery distributed worker nodes orchestrated via RabbitMQ message brokers. \\
  \addlinespace
  \textbf{Model Inference} & High memory consumption if complex tree ensembles are scaled. & Retain the compiled Weight-of-Evidence lookup table scorecard in memory; C-level lookup executes in nanoseconds without tree traversal. \\
  \bottomrule
\end{tabularx}
\caption{Technical Bottleneck Identification and Production Engineering Mitigations.}
\label{tab:bottleneck_analysis}
\end{table}
